\section{Background and Related Work}
\label{sec:related}

\paragraph{Probing and its controls.}
Linear classifiers have long been used to track information across network depth \citep{alain2017understanding}. A probe score reflects both the representation and the learner fitted to it. This motivates matched-capacity probes and control tasks \citep{hewitt2019designing}. The probing literature distinguishes information a readout can recover from information the model uses \citep{belinkov2022probing,ravichander2021probing}. Causal erasure pursues the complementary test \citep{ravfogel2020null,elazar2021amnesic}, and a closed-form projection removes a concept from every linear reader at once \citep{belrose2023leace}. Erasure asks what breaks when a concept is deleted; the question here is whether the concept's own direction is the one that moves its answer. Our readability grade adopts the control-task idea: it measures selectivity against random-label probes of equal capacity fitted within recurring image types. Depth-wise probing of medical multimodal models has mapped classification information across towers, connectors, and language layers \citep{zhu2026lost}. Layer-wise probes of MedCLIP expose acquisition and device shortcuts \citep{pedersen2026probes}. Neither intervenes on the decoded clinical direction at the same locus.

\paragraph{Concept directions and what they inherit.}
A concept activation vector is the normal of a classifier separating concept examples from a negative set, and it grades a class by the directional derivative of the logit along that normal \citep{kim2018tcav}. We keep the direction and change the measurement: a bounded forward write at a fixed fraction of each token's norm, scored against directions built the same way for the competing concepts, rather than a gradient sensitivity read at the output. Such a normal shifts with the layer it is fitted at, entangles with correlated concepts, depends on the negative set it is fitted against, and carries where in the image the concept was read \citep{nicolson2025explaining}. Section~\ref{sec:robust-geometry} tests each: alternative constructions, a second locus, and writes concentrated on the tokens the probe reads.

\paragraph{Steering with probe normals.}
Representation engineering makes the step from reading a concept to writing it back \citep{zou2023representation}. Activation engineering changes intermediate representations at inference time and measures the resulting output changes \citep{turner2023steering,panickssery2024caa}. A probe normal is an appealing intervention because it supplies a label and a signed vector from the same representation. Separability-trained concept vectors can absorb distractor signals and diverge from the concept they are meant to model \citep{pahde2025navigating}. A classifier's weight vector is a filter, not the pattern it detects \citep{haufe2014interpretation}. Mass-mean directions steer truth judgements in language models more effectively than logistic normals \citep{marks2024geometry}. The features that identify a concept in the input are largely not the features that control the output, and selecting for output effect improves steering \citep{arad2025saes}. Benchmarks accordingly score concept detection and steering as separate axes of the same method \citep{wu2025axbench}. Random-direction and sham controls test generic sensitivity and coordinate dependence. They do not test whether an equally constructed direction with a different label would move the same answer more. Ownership adds this comparison and evaluates it in the forward path at the consumed locus, using each token's norm as the dose scale. Alignment search certifies a subspace as causal by showing that it implements an interpretable causal variable \citep{geiger2024alignments}; ownership asks a narrower question of a labelled direction: whether the label names the direction that moves its answer. Prior work replaces the separability-trained vector with a better estimator of the same concept and evaluates the replacement against the concept's own label \citep{pahde2025navigating,haufe2014interpretation,marks2024geometry}. Ownership instead fixes the estimator and varies the label: the reference family contains a direction fitted the same way for every other concept in the set, so a construction that improves estimation is scored by whether it wins that comparison. Section~\ref{sec:robust-geometry} runs five constructions and both lifts through the same test. Geometric notions such as cosine and projection depend on the inner product chosen for the representation \citep{park2024linear}, so every cosine here names its metric: cosines between directions are reported in the raw model space, and the two comparisons that carry the argument---the answer direction against the probe normal, and the expert-label refit against the report-label direction---in both the raw model space and the training covariance. It scores difference-of-means, pattern, orthogonalised, and residualised directions alongside the logistic normal (Section~\ref{sec:robust-geometry}).

\paragraph{Writing into a vision-language model.}
Image features are linearly decodable in a vision-language model's hidden layers, and targeted edits to them change downstream answers \citep{rajaram2025lineofsight}. Steering one sparse-autoencoder feature of the vision tower on an empty image and asking the language model what it sees turns the language model into a read-out of what that feature encodes \citep{ferrando2026explain}. That work steers a feature to recover its meaning; this paper starts from the meaning, fits a direction to a clinical label, and asks whether the write keeps that label's selectivity when the other findings' directions are written at the same locus, dose, and patients. Causal probing of multimodal models intervenes on the language model's residual stream with difference-in-means vectors over paired images, and reads out where entity and abstract concepts sit across depth \citep{deng2026causal}. It intervenes inside the reader and scores each concept against its own effect; the write here acts on the representation the reader consumes and scores each concept against the competing concepts' directions, which is the comparison that decides whether the label names the handle.
